\chapter{Трёхузловой parent graph: спектральный rank selector без ручного знака}

Предыдущие gates нашли правильный wedge-channel, но оставили открытым
происхождение relative minus между direct и crossed paths. Здесь проверяется
более экономная конструкция: не добавлять разность каналов вручную, а провести
поле \(\Delta\) через внешний квадрат. Оказывается, полный raw spectral trace
такого графа сам создаёт положительный determinant invariant.

\section{Структура предлагаемого механизма}

Предыдущий анализ определил требуемый determinant stabilizer, но не его
parent construction. Рассмотрим трёхузловую цепь:
\begin{equation}
 \mathbb C
 \xrightarrow{\ A_\Delta\ }
 \mathbb C^2\otimes\mathbb C^4
 \xrightarrow{\ B_\Delta\ }
 \Lambda^2\mathbb C^2\otimes\Lambda^2\mathbb C^4.
 \label{eq:ps-three-node-chain}
\end{equation}
Первое ребро кладёт \(\Delta\) в восьмимерный intermediate space. Второе
ребро берёт все его \(2\times2\) minors. Поэтому двухшаговый путь видит ровно
то, превращается ли rank one в rank two. Никакой отдельный знак
\(H^{\rm dir}-H^{\rm cr}\) больше не вставляется: antisymmetrization уже
содержится в \(\Lambda^2\).

\section{Два фундаментальных ребра}

Положим
\begin{equation}
 A_\Delta(1)=\operatorname{vec}\Delta.
\end{equation}
Для \(x\in\mathbb C^2\otimes\mathbb C^4\) определим
\begin{equation}
 (B_\Delta x)_{IJ}
 =\frac{c}{2}\left(
 \Delta_{0I}x_{1J}-\Delta_{1I}x_{0J}
 -\Delta_{0J}x_{1I}+\Delta_{1J}x_{0I}
 \right),
 \qquad I<J.
 \label{eq:ps-second-wedge-edge}
\end{equation}
Здесь \(c\) временно сохраняет возможную относительную нормировку двух
representations одной one-form. При канонической polarization convention
\(c=1\). Композиция равна
\begin{equation}
 \boxed{B_\Delta A_\Delta=c\,\Lambda^2\Delta},
 \qquad
 \|\Lambda^2\Delta\|^2
 =\det(\Delta\Delta^\dagger).
 \label{eq:ps-three-node-minor-identity}
\end{equation}

\section{Raw spectral traces}

На graded carrier размерности \(1+8+6\) возьмём
\begin{equation}
 D_\Delta=
 \begin{pmatrix}
 0&A_\Delta^\dagger&0\\
 A_\Delta&0&B_\Delta^\dagger\\
 0&B_\Delta&0
 \end{pmatrix}.
 \label{eq:ps-three-node-dirac}
\end{equation}
Обозначим
\begin{equation}
 \rho=\Tr(\Delta\Delta^\dagger),\qquad
 \tau=\Tr[(\Delta\Delta^\dagger)^2],\qquad
 d=\det(\Delta\Delta^\dagger).
\end{equation}
Прямое умножение даёт
\begin{align}
 \Tr D_\Delta^2
 &=\left(2+\frac32c^2\right)\rho,\\
 \Tr D_\Delta^4
 &=2\rho^2+c^4\left(\frac38\tau+\frac14d\right)
   +4c^2d\\
 &=\left(2+\frac38c^4\right)\rho^2
   +\left(4c^2-\frac12c^4\right)d.
 \label{eq:ps-three-node-traces}
\end{align}
Последний переход использует \(\tau=\rho^2-2d\), справедливое для
\(2\times4\) matrix. Следовательно, determinant coefficient положителен не
в одной подогнанной точке, а во всём интервале
\begin{equation}
 \boxed{0<c^2<8.}
 \label{eq:ps-three-node-stability-window}
\end{equation}

При \(c=1\) physical half-traces KO6 completion равны
\begin{equation}
 \frac12\Tr D_F^2=\frac72\rho,
 \qquad
 \frac12\Tr D_F^4=\frac{19}{8}\rho^2+\frac72d.
 \label{eq:ps-three-node-canonical-traces}
\end{equation}
То есть стабилизатор появляется из обычного незавешенного spectral trace, а
не из Casimir-decorated functional.

\section{Vacuum и полный Hessian}

Для нормировки
\begin{equation}
 V=-\frac12\Tr D_F^2+\frac12\Tr D_F^4
 =-\frac72\rho+\frac{19}{8}\rho^2+\frac72d
 \label{eq:ps-three-node-potential}
\end{equation}
rank-one minimum имеет
\begin{equation}
 \rho_*=\frac{14}{19},
 \qquad
 V_*=-\frac{49}{38}.
\end{equation}
Equal-singular-value rank-two branch имеет
\begin{equation}
 p=q=\frac7{26},
 \qquad
 V_{(2)}=-\frac{49}{52}>V_*.
\end{equation}
Exact real \(16\times16\) Hessian spectrum в rank-one точке равен
\begin{equation}
 \boxed{
 14\ (1),\qquad
 0\ (9),\qquad
 \frac{98}{19}\ (6).}
 \label{eq:ps-three-node-hessian}
\end{equation}
Таким образом, radial mode положителен, девять gauge-orbit directions
остаются Goldstone, а все шесть rank-growth directions становятся
положительными.

\section{KO6 и Krajewski-проверка}

Particle chain имеет grading \((+,-,+)\). Добавим conjugate chain и зададим
\(J_F\) обменом двух copies с complex conjugation. На complex dimension 30
machine audit проверяет
\begin{equation}
 D_F=D_F^\dagger,\qquad
 \{D_F,\Gamma_F\}=0,\qquad
 [D_F,J_F]=0,\qquad
 \{J_F,\Gamma_F\}=0
\end{equation}
с ошибками порядка machine precision.

Coarse labels можно формально выбрать как
\begin{equation}
 (0,0)\longrightarrow(0,1)\longrightarrow(1,1),
\end{equation}
а для conjugate chain --- transposed labels. Каждое fundamental edge меняет
ровно один label, поэтому chain проходит необходимый coarse order-one edge
test. Однако это ещё не доказывает существование самих трёх nodes как
modules represented associative algebra. Следующий audit показывает, что
для \(\mathcal A_F^{PS}=\mathbb H_R\oplus\mathbb H_L\oplus M_4(\mathbb C)\)
одномерный node и irreducible color-six node не существуют. Поэтому coarse
labeling нельзя повышать до настоящего Krajewski pass.

\begin{theorem}[Equivariant carrier pass]
Трёхузловая exterior-square chain имеет KO6 completion и проходит coarse
order-one edge criterion. Её polynomial trace автоматически содержит
положительный rank-sensitive determinant coefficient и выбирает требуемый
rank-one orbit с Hessian signature \((7_+,9_0,0_-)\). Relative
direct-minus-crossed sign фиксирован antisymmetrization внешнего квадрата.
Но этот результат является graded superconnection carrier, а не строгим
finite spectral triple над неизменённой Pati--Salam algebra.
\end{theorem}

\section{Что ещё не закрыто}

Это gauge-equivariant parent carrier, но не finite-triple parent graph. В
\eqref{eq:ps-second-wedge-edge} оба ребра связаны с одной и той же
\(\Delta\) по invariant polarization prescription, а endpoint
\((1_R,6_4)\) является допустимым scalar channel, но не module node
\(M_4(\mathbb C)\). Поэтому literal three-node embedding закрывается.
Рабочий механизм следует искать в represented universal two-forms или
curvature внутри уже существующего 32-dimensional Pati--Salam Hilbert space,
после quotient по degree-two junk.

\section{Литература и воспроизводимость}

Diagrammatic finite-triple interpretation сверена с T.~Krajewski,
\emph{Classification of Finite Spectral Triples}, arXiv:hep-th/9701081.
Связь finite geometry со scalar-potential constraints сверена с
T.~Krajewski, \emph{Constraints on the Scalar Potential from Spectral Action
Principle}, arXiv:hep-th/9803199. Использование обычных следов
\(\Tr D_F^2\) и \(\Tr D_F^4\) опирается на A.~H.~Chamseddine and A.~Connes,
\emph{The Spectral Action Principle}, arXiv:hep-th/9606001.

Machine audit сохранён в
\path{s2t_v4_pati_salam_three_node_parent_graph.py} и
\path{s2t_v4_pati_salam_three_node_parent_graph_results.json}.